\clearpage
\section{Предварительная обработка данных}

Диапазон изменения входных данных составляет до 6 порядков -- от -60~дБ до +60~дБ. При этом, изменения характеристик в отдельных участках частот могут составлять единицы децибел. В связи с этим, типовые нормировки вида norm/std, min/max делают модель нечувствительной к малым изменениям характеристики.

Принято решение перейти от абсолютных значений амплитуды и фазы к наклонам этих характеристик. Наклоны представлены в ${C_\text{in}=4}$ каналах. Длина каждого канала приводится к ${L_1=512}$ путем интерполяции. Массивы должны соответствовать логарифмическому масштабу по частоте. В примере на рисунках~\labelcref{fig:preproc,fig:preproc_augmented} на последних 4 графиках изображены данные, подаваемые на вход модели сегментации.

\clearpage
\begin{figure}[H]
	\centering
	\includegraphics [width=1.00\linewidth] {parametersEstimation/dataset_samples.pdf}

	\caption{Пример разбиения данных на 4 канала}
	\label{fig:preproc}
\end{figure}

\clearpage
\begin{figure}[H]
	\centering
	\includegraphics [width=1.00\linewidth] {parametersEstimation/augmented_dataset_samples.pdf}

	\caption{Пример разбиения данных на 4 канала, с аугментациями}
	\label{fig:preproc_augmented}
\end{figure}
